\scanpage{49}
- 211 - \hfill n\textsuperscript{o} 391

(III) En vertu du § 8, n\textsuperscript{o} 3, formule (4), on a $h=n/\ell=30$.

(IV) La plus grande racine pour l'ordre opposé de l'ordre lexicographique sur $\mathrm{R}^8$ est évidemment $\varepsilon_7+\varepsilon_8=2\alpha_1+3\alpha_2+4\alpha_3+6\alpha_4+5\alpha_5+4\alpha_6+3\alpha_7+2\alpha_8$ ; c'est la plus grande racine $\widetilde\alpha$, (§ 6, n\textsuperscript{o} 6, cor. 2 de la prop. 13), et on vérifie que son ordre est $h-1=29$. Elle est orthogonale à toutes les racines simples sauf $\alpha_8$, et l'on a $(\widetilde\alpha|\alpha_8)=-1$.

(V) Les \uline{coracines} s'obtiennent en remplaçant les $\varepsilon_i$ par les $\varepsilon_i^*$ dans les expressions des racines de $R$, et $R^*$ est donc isomorphe à $R$. Comme toutes les racines sont de même longueur, la \uline{forme bilinéaire canonique} pour $R$ est $(x|y)/2h=(x|y)/60$ et la longueur des racines est $1/\sqrt{30}$, (§ 8, n\textsuperscript{o} 5, \uline{Remarque}). Le nombre des racines non orthogonales à une racine donnée est $4h-6=114$ et le nombre $\gamma$ est égal à $h^2=900$.

(VI) Comme $\Gamma_{\unsure{o|0}}(R^*)$ est le réseau $L_3(V^*)$, son associé dans $V=\mathrm{R}^4$ est $L_3(V)$, (n\textsuperscript{o} 4), c'est-à-dire $\Gamma_{\unsure{o|0}}(R)$, autrement dit, on a $\Gamma(R)=\Gamma_{\unsure{o|0}}(R)$. Le calcul des poids fondamentaux, formant \unsure{le|la} base duale de $(\alpha_i^*)$, donne :
\[
\omega_1=2\varepsilon_8=4\alpha_1+5\alpha_2+7\alpha_3+10\alpha_4+8\alpha_5+6\alpha_6+4\alpha_7+2\alpha_8
\]
\[
\begin{aligned}
\omega_2&=\frac12(\varepsilon_1+\varepsilon_2+\varepsilon_3+\varepsilon_4+\varepsilon_5+\varepsilon_6+\varepsilon_7+5\varepsilon_8)=\\
&=5\alpha_1+8\alpha_2+10\alpha_3+15\alpha_4+12\alpha_5+9\alpha_6+6\alpha_7+3\alpha_8
\end{aligned}
\]
\[
\begin{aligned}
\omega_3&=\frac12(-\varepsilon_1+\varepsilon_2+\varepsilon_3+\varepsilon_4+\varepsilon_5+\varepsilon_6+\varepsilon_7+7\varepsilon_8)=\\
&=7\alpha_1+10\alpha_2+14\alpha_3+20\alpha_4+16\alpha_5+12\alpha_6+8\alpha_7+4\alpha_8
\end{aligned}
\]
\[
\begin{aligned}
\omega_4&=\varepsilon_3+\varepsilon_4+\varepsilon_5+\varepsilon_6+\varepsilon_7+5\varepsilon_8=\\
&=10\alpha_1+15\alpha_2+20\alpha_3+30\alpha_4+24\alpha_5+18\alpha_6+12\alpha_7+6\alpha_8
\end{aligned}
\]
\mspage{48}
